\documentclass[10pt]{letter}
\usepackage[T1]{fontenc}
\usepackage[utf8]{inputenc}
\usepackage{lmodern}
\usepackage{microtype}
\usepackage[a4paper,margin=2.35cm]{geometry}
% Author metadata supplied by the author team on 2026-08-11.
\newif\ifAuthorMetadataReady
\AuthorMetadataReadytrue

\newcommand{\CorrespondingAuthor}{Zhuo Chen}
\newcommand{\CorrespondingEmail}{zhuoc@chalmers.se}
\newcommand{\CorrespondingAffiliation}{Chalmers University of Technology, SE-412 96 Gothenburg, Sweden}

\newcommand{\AuthorFrontmatter}{%
  \def\thefnote{$\dagger$}%
  \author[aff-a]{Shuhao Liu\fnref{equal}}
  \ead{2420610137@stu.hrbust.edu.cn}
  \author[aff-b]{Zhuo Chen\fnref{equal}\corref{corresponding}}
  \ead{zhuoc@chalmers.se}
  \author[aff-c]{Zhi Ling}
  \ead{zhi.ling@kiwiar.com}
  \author[aff-c]{Yu Yan}
  \ead{yu.yan@kiwiar.com}
  \author[aff-a]{Jie Liu}
  \ead{liujie@hrbust.edu.cn}
  \author[aff-a]{Qiuxue Wu}
  \ead{2420610140@stu.hrbust.edu.cn}
  \author[aff-c]{Ziyi Kuang}
  \ead{ziyi.kuang@kiwiar.com}
  \address[aff-a]{Harbin University of Science and Technology, No. 52 Xuefu Road, Nangang District, Harbin 150080, Heilongjiang, China}
  \address[aff-b]{Chalmers University of Technology, SE-412 96 Gothenburg, Sweden}
  \address[aff-c]{Kiwiar Co., Ltd., Singapore}
  \fntext[equal]{These authors contributed equally to this work.}
  \cortext[corresponding]{Corresponding author.}
}

% These statements require separate confirmation from the author team.
\newif\ifFundingStatementReady
\FundingStatementReadyfalse
\newcommand{\FundingStatement}{}
\newif\ifCompetingInterestsReady
\CompetingInterestsReadyfalse
\newcommand{\CompetingInterests}{}
\newif\ifCRediTStatementReady
\CRediTStatementReadyfalse
\newcommand{\CRediTStatement}{}
\newcommand{\RepositoryDOI}{}

\ifAuthorMetadataReady
\signature{\CorrespondingAuthor\\on behalf of all authors}
\address{\CorrespondingAffiliation\\Email: \CorrespondingEmail}
\else
\signature{Corresponding author\\on behalf of all authors}
\fi
\begin{document}
\begin{letter}{Editors-in-Chief\\Smart Agricultural Technology}
\opening{Dear Editors,}

Please consider the manuscript, ``Support-Aware Review Queues for Crop--Weed Vision: A Multitemporal Single-Field Case Study,'' for publication in \emph{Smart Agricultural Technology}.

The study addresses a practical commissioning question for agricultural vision: which stage currently limits a crop--weed review queue, candidate generation, role qualification, rank ordering, or allocation of a finite review budget? We answer this question on the official WeedsGalore spatial split using instance-level denominators from proposal formation through final review selection.

The framework separates spatial proposal recall, role-qualified recall, and queued recall under many-to-one and one-to-one matching. Shared-ranker-weight comparisons isolate proposal formation, while all-candidate fitting aligns ranker supervision with queue support. Candidate-level role-qualified AUC and average precision are linked to instance recovery, crop exposure, and the finite review budget rather than treated as end-to-end system metrics.

A target-semantic baseline with equal target-mask access further distinguishes component formation from ranking. Across five independent H200 training seeds, fixed-$K$ precision and one-to-one recall were $0.7168\pm0.0260$ and $0.2027\pm0.0057$. Exploratory 5--50 candidates/tile allocation improved both metrics in every seed, to $0.8939\pm0.0249$ and $0.2578\pm0.0090$, while worst-date recall decreased in every seed. A complete 16-policy quota grid and date-balanced comparator show how pooled utility, weakest-date coverage, and spatial-block effects diverge. Validation selection retained fixed $K=20$ in four seeds and 0--30 allocation in one.

The manuscript presents an application-focused queue evaluation as a multitemporal single-field case study. Its chronology distinguishes the archived core analysis from post-test exploratory extensions. The accompanying package includes editable figures, source-data tables, exact component-mask matching outputs, complete proposal and quota grids, versioned configurations, failure records, run hashes, runtime and memory records, and clean-room replay commands.

\closing{Sincerely,}
\end{letter}
\end{document}
